% GENERATED FILE. Edit canonical lessons, then run npm run generate:books.
% canonical-source-hash: fnv1a64:140d6e26b2246da4
% canonical-lessons: ZH-C524-jiezhang, ZH-C524-banjia, ZH-W524-jia-price, ZH-C524-yaojia, ZH-C524-huanjia, ZH-C524-wujia

\chapter{Half Price}
\label{ch:zh-half-price}
\hlchaptermodality{524}

\begin{chapteropening}
\textbf{By the end of this chapter:} \emph{I can talk about paying the bill, half price, the asking price, offering less and the cost of living.}

Complete the last lesson of chapter 524: I can talk about paying the bill, half price, the asking price, offering less and the cost of living.
\end{chapteropening}

\section[jiézhàng]{jiézhàng — to settle the bill}

\label{lesson:ZH-C524-jiezhang}



\begin{quote}
Before the new one: say the Mandarin for the buyer (in a sale), then the Mandarin for to sell retail.
\end{quote}

\subsection*{You'll want to know: jiézhàng}
\emph{jiézhàng} — \textquotedblleft{}to settle the bill\textquotedblright{}.

\begin{quote}
\emph{wǒmen děng le èrshí fēnzhōng cái jiézhàng} — We waited twenty minutes before we could settle the bill.
\end{quote}

Its characters come later; for now it is a word for the ear.

\begin{cousinweb}
\emph{jiézhàng} is built from \emph{jié} 'to tie' and \emph{zhàng} 'an account'. The first is inherited from Middle Chinese \emph{ket} (Guangyun, in Baxter's transcription), glossed 'to tie'; no dictionary we use records an older form of the second, so its history is not settled.
\end{cousinweb}

\begin{grammarlens}[title={What you've built}]
One more word for saying how it was.
\end{grammarlens}

\subsection*{Your turn}
\begin{itemize}
\raggedright
  \item \emph{Say it:} \emph{jiézhàng}
  \item \emph{Say it:} \emph{jiézhàng}, once more
  \item \emph{Say it:} \emph{wǒmen děng le èrshí fēnzhōng cái jiézhàng} — We waited twenty minutes before we could settle the bill.
  \item \emph{Recall:} say the Mandarin for the buyer (in a sale), then the Mandarin for to sell retail, then say \emph{jiézhàng} again
\end{itemize}

\begin{tcolorbox}[breakable,colback=teal!4,colframe=teal!35!black,title={Before you move on}]
What does jiézhàng mean? (\textquotedblleft{}To settle the bill\textquotedblright{}.) Say it once more.
\end{tcolorbox}
\section[\texorpdfstring{\zh{半价} (bànjià)}{半价 (bànjià)}]{\zh{半价} (bànjià) — half price}

\label{lesson:ZH-C524-banjia}



\begin{quote}
Before the new one: say the Mandarin for to sell retail, then the Mandarin for to settle the bill.
\end{quote}

\subsection*{You'll want to know: \zh{半价}}
\textbf{\zh{半价}} — \emph{bànjià} — \textquotedblleft{}half price\textquotedblright{}.

\begin{quote}
\textbf{\zh{今天这儿的东西都半价}} \emph{jīntiān zhèr de dōngxi dōu bànjià} — Everything here is half price today.
\end{quote}

The new character, in pieces: \zh{价} — six strokes in three pieces: the person piece on the left, two strokes, as in \zh{你}; \zh{人} on the right, two strokes; and under \zh{人} a short stroke falling to the left and a vertical. The next lesson puts it in your hand.

\begin{cousinweb}
\textbf{\zh{半价}} (\emph{bànjià}) is built from \zh{半} 'half' and \zh{价} 'a price', the jià of the word for price, both inherited from Middle Chinese, \emph{panH} and \emph{kaeH} (Guangyun, in Baxter's transcription), which glosses the first as a thing split in the middle and the second as the count of a price.
\end{cousinweb}

\begin{grammarlens}[title={What you've built}]
One more word for saying how it was.
\end{grammarlens}

\subsection*{Your turn}
\begin{itemize}
\raggedright
  \item \emph{Say it:} \emph{bànjià}
  \item \emph{Say it:} \emph{bànjià}, once more
  \item \emph{Say it:} \emph{jīntiān zhèr de dōngxi dōu bànjià} — Everything here is half price today.
  \item \emph{Recall:} say the Mandarin for to sell retail, then the Mandarin for to settle the bill, then say \emph{bànjià} again
\end{itemize}

\begin{tcolorbox}[breakable,colback=teal!4,colframe=teal!35!black,title={Before you move on}]
What does \zh{半价} mean? (\textquotedblleft{}Half price\textquotedblright{}.) Say it once more.
\end{tcolorbox}
\section[\texorpdfstring{\zh{价} (jià)}{价 (jià)}]{\zh{价} — the person piece, then \zh{人} with two legs under it}

\label{lesson:ZH-W524-jia-price}



\begin{quote}
The word first, then the character inside it.

\begin{itemize}
\raggedright
  \item \emph{Recall:} say \emph{bànjià} — \textbf{R1}, one lesson back, before it decays
\end{itemize}

\begin{quote}
\zh{你} \zh{人}
\end{quote}

\emph{Read it:} each one aloud, with nothing said first — \emph{nǐ}, \emph{rén}
\end{quote}

\subsection*{Script — the person piece, then \zh{人} with two legs under it}
\hlblockfigure{\detokenize{figures/ZH-W524-jia-price-filmstrip.pdf}}{How \zh{价} is written, stroke by stroke}

\begin{quote}
\zh{价}
\end{quote}

\textbf{Three pieces: the person piece, then \zh{人}, then what stands under it.} Count the pieces, not the strokes.

\textbf{The left piece, the person piece, two strokes}, as in \zh{你}: a stroke falling to the left, then a vertical down to the foot.

\textbf{The top right piece, \zh{人}, two strokes}, as you wrote it, spread wide like a roof: a long stroke falling to the left, then a long stroke falling to the right.

\textbf{Under it, two strokes}: a short stroke that drops and falls to the left, then a vertical on the right.

2 + 2 + 2 = six.

\subsection*{Your turn}
\begin{itemize}
\raggedright
  \item \emph{Write it:} \zh{价}, all six, with the model in front of you
  \item \emph{Cover:} write it again from memory, and count the pieces, not the strokes
\end{itemize}

\begin{tcolorbox}[breakable,colback=teal!4,colframe=teal!35!black,title={Before you move on}]
How many strokes does \zh{价} have? (\textbf{Six}.) What is the first piece? (\textbf{The person piece, on the left}.)

Source: \href{https://raw.githubusercontent.com/chanind/hanzi-writer-data/68d10a4b21150cae5e1ebbd223eed289cf32d90c/data/\%E4\%BB\%B7.json}{Hanzi Writer Data \zh{价}, pinned stroke paths and medians}.
\end{tcolorbox}
\section[\texorpdfstring{\zh{要价} (yàojià)}{要价 (yàojià)}]{\zh{要价} (yàojià) — to name one's price; to ask a price}

\label{lesson:ZH-C524-yaojia}



\begin{quote}
Before the new one: say the Mandarin for to settle the bill, then the Mandarin for at half price.
\end{quote}

\subsection*{You'll want to know: \zh{要价}}
\textbf{\zh{要价}} — \emph{yàojià} — \textquotedblleft{}to name one's price; to ask a price\textquotedblright{}.

\begin{quote}
\emph{tā yàojià bǐ nà jiā shāngdiàn gāo} — He's asking more than that shop.
\end{quote}

\begin{cousinweb}
\textbf{\zh{要价}} (\emph{yàojià}) is built from \zh{要} 'to want' and \zh{价} 'a price', both inherited from Middle Chinese, \emph{'jiewH} and \emph{kaeH} (Guangyun, in Baxter's transcription), which glosses the first 'to bind by agreement' and the second as the count of a price.
\end{cousinweb}

\begin{grammarlens}[title={What you've built}]
One more word for saying how it was.
\end{grammarlens}

\subsection*{Your turn}
\begin{itemize}
\raggedright
  \item \emph{Say it:} \emph{yàojià}
  \item \emph{Say it:} \emph{yàojià}, once more
  \item \emph{Say it:} \emph{tā yàojià bǐ nà jiā shāngdiàn gāo} — He's asking more than that shop.
  \item \emph{Recall:} say the Mandarin for to settle the bill, then the Mandarin for at half price, then say \emph{yàojià} again
\end{itemize}

\begin{tcolorbox}[breakable,colback=teal!4,colframe=teal!35!black,title={Before you move on}]
What does \zh{要价} mean? (\textquotedblleft{}To name one's price; to ask a price\textquotedblright{}.) Say it once more.
\end{tcolorbox}
\section[\texorpdfstring{\zh{还价} (huánjià)}{还价 (huánjià)}]{\zh{还价} (huánjià) — to offer a lower price back; to make a counter-offer}

\label{lesson:ZH-C524-huanjia}



\begin{quote}
Before the new one: say the Mandarin for at half price, then the Mandarin for to name one's price.
\end{quote}

\subsection*{You'll want to know: \zh{还价}}
\textbf{\zh{还价}} — \emph{huánjià} — \textquotedblleft{}to offer a lower price back; to make a counter-offer\textquotedblright{}.

\begin{quote}
\textbf{\zh{在这儿买东西最好还价}} \emph{zài zhèr mǎi dōngxi zuìhǎo huánjià} — When you buy things here, it's best to offer less.
\end{quote}

\begin{cousinweb}
\textbf{\zh{还价}} (\emph{huánjià}) is built from \zh{还} huán 'to give back' and \zh{价} 'a price', both inherited from Middle Chinese, \emph{hwaen} and \emph{kaeH} (Guangyun, in Baxter's transcription), which glosses the first 'to return': a price handed back.
\end{cousinweb}

\begin{grammarlens}[title={What you've built}]
One more word for saying how it was.
\end{grammarlens}

\subsection*{Your turn}
\begin{itemize}
\raggedright
  \item \emph{Say it:} \emph{huánjià}
  \item \emph{Say it:} \emph{huánjià}, once more
  \item \emph{Say it:} \emph{zài zhèr mǎi dōngxi zuìhǎo huánjià} — When you buy things here, it's best to offer less.
  \item \emph{Recall:} say the Mandarin for at half price, then the Mandarin for to name one's price, then say \emph{huánjià} again
\end{itemize}

\begin{tcolorbox}[breakable,colback=teal!4,colframe=teal!35!black,title={Before you move on}]
What does \zh{还价} mean? (\textquotedblleft{}To offer a lower price back; to make a counter-offer\textquotedblright{}.) Say it once more.
\end{tcolorbox}
\section[wùjià]{wùjià — prices; the cost of living}

\label{lesson:ZH-C524-wujia}



\begin{quote}
Before the new one: say the Mandarin for to name one's price, then the Mandarin for to offer a lower price back.
\end{quote}

\subsection*{You'll want to know: wùjià}
\emph{wùjià} — \textquotedblleft{}prices; the cost of living\textquotedblright{}.

\begin{quote}
\emph{zhèr de wùjià bǐ Běijīng gāo} — The cost of living here is higher than in Beijing.
\end{quote}

Its characters come later; for now it is a word for the ear.

\begin{cousinweb}
\emph{wùjià} is built from \emph{wù} 'a thing' and \zh{价} 'a price', both inherited from Middle Chinese, \emph{mjut} and \emph{kaeH} (Guangyun, in Baxter's transcription), which glosses the first 'all things': the price of things.
\end{cousinweb}

\begin{grammarlens}[title={What you've built}]
One more word for saying how it was.
\end{grammarlens}

\subsection*{Your turn}
\begin{itemize}
\raggedright
  \item \emph{Say it:} \emph{wùjià}
  \item \emph{Say it:} \emph{wùjià}, once more
  \item \emph{Say it:} \emph{zhèr de wùjià bǐ Běijīng gāo} — The cost of living here is higher than in Beijing.
  \item \emph{Recall:} say the Mandarin for to name one's price, then the Mandarin for to offer a lower price back, then say \emph{wùjià} again
\end{itemize}

\begin{tcolorbox}[breakable,colback=teal!4,colframe=teal!35!black,title={Before you move on}]
What does wùjià mean? (\textquotedblleft{}Prices; the cost of living\textquotedblright{}.) Say it once more.
\end{tcolorbox}
